% P06.2 skeletal-current repair, exact round-2 packet, 25 August 2026.
% The rejected round-1 boundary-separation proof is not used.
% No novelty or priority claim is intended.

\paragraph{Source scope.}
The proof uses three primary sources.
\begin{enumerate}
\item In Gubler--Jell--Rabinoff, \emph{Forms on Berkovich spaces based on
harmonic tropicalizations}, arXiv:2111.05741v3, Proposition~5.8 and
Lemma~5.10 give a local decomposition of an $\mathbf R$-PL function into
integral PL functions with $\mathbf Q$-linearly independent real
coefficients.  Proposition~8.3 and paragraph~8.4 give a common smooth
tropicalization for a finite tuple of integral PL functions.
\item In Chambert--Loir--Ducros, \emph{Formes diff\'erentielles r\'eelles et
courants sur les espaces de Berkovich}, arXiv:1204.6277v2 (25 July 2025),
Definition~7.5.1 defines skeletal currents by calibrated skeleta in every
tropical chart; Remark~7.2.5 says that the localized chart tests occurring
in that definition determine a current; Example~7.5.3 says that the
integration current of a boundaryless analytic space is skeletal with its
canonical calibration; Propositions~7.5.7 and~7.5.8 give compatibility and
uniqueness of calibrated skeletal descriptions under common refinement;
and Theorem~8.10.3 says that $\mathrm d'\mathrm d''(uT)$ is skeletal when
$T$ is closed skeletal and $u$ is PL.  More precisely, the proof of that
theorem computes the calibration along a codimension-one cell as the sum of
the normal slopes of $u$ on the adjacent top-dimensional cells, weighted
by their calibrations.  Theorem~18.8.4 with $d=1$ identifies vanishing of
the curvature current of an integral PL metric with numerical triviality
of its residual first Chern class.
\item In Gubler--Rabinoff, \emph{Classically psh and pluriharmonic
functions on Berkovich spaces}, arXiv:2506.13548v2, Theorem~4.3 identifies
semipositivity of an $\mathbf R$-PL function with classical
plurisubharmonicity over a nontrivially valued field.
\end{enumerate}
The audited PDF SHA-256 values, in the same order, are
\texttt{572d51c9f8907dbb814684f7d136612243c0e0b70cddbe6759bb686f0bcfbd42},
\texttt{5690f999523017d184ab751bb55857ab1637b197d1186ad4fd2edba7e1bda719},
and
\texttt{83ca9a19394f44bb2e4e1aa316de95938f311acd53a00fe250fa0f605184b726}.

\begin{proposition}
Let $K$ be a complete nontrivially valued non-Archimedean field, and let
$X$ be a good, separated, boundaryless, strictly $K$-analytic space that is
pure of dimension $n\geq1$.  If $h:X\to\mathbf R$ is an $\mathbf R$-PL
function and its Chambert--Loir--Ducros current satisfies
\[
                         \mathrm d'\mathrm d''[h]=0,
\]
then $h$ is harmonic in the sense of Gubler--Jell--Rabinoff.  Equivalently,
the residue class $L_h(x)\in\operatorname{Pic}(\operatorname{red}(X,x))_
{\mathbf R}$ is numerically trivial for every $x\in X$.  Consequently $h$
is classically pluriharmonic in the sense of Gubler--Rabinoff.
\end{proposition}

\begin{proof}
Fix $x\in X$.  By Gubler--Jell--Rabinoff, Proposition~5.8 and
Lemma~5.10, after shrinking to an analytic neighbourhood $W$ of $x$, one
can write
\begin{equation}
\label{eq:p06r2-independent}
                         h=\sum_{i=1}^{s}\lambda_i h_i,
\end{equation}
where the $h_i$ are integral PL functions on $W$ and the real numbers
$\lambda_1,\ldots,\lambda_s$ are linearly independent over $\mathbf Q$.
We prove that
\begin{equation}
\label{eq:p06r2-integral-currents}
                         \mathrm d'\mathrm d''[h_i]=0
                         \qquad(1\leq i\leq s)
\end{equation}
on a neighbourhood of $x$.

Let $T_0=[W]$ be the integration current.  All assertions are local, so we
may replace $W$ by a smaller open subset and suppose it is boundaryless and
pure of dimension $n$.  By CLD Example~7.5.3, $T_0$ is a closed skeletal
current of dimension $n$, described in each moment map by the canonical
tropical calibration.  For each $i$, set
\[
                         T_i:=\mathrm d'\mathrm d''(h_iT_0)
                             =\mathrm d'\mathrm d''[h_i].
\]
CLD Theorem~8.10.3 says that $T_i$ is a skeletal current of dimension
$n-1$.

We recall the arithmetic in the construction in the proof of that theorem.
Choose a smooth tropicalization $F$ on an open subset $U\subset W$ which
simultaneously tropicalizes $h_1,\ldots,h_s$.  Thus
$h_i=\varphi_i\circ F$, with $\varphi_i$ piecewise
$(\mathbf Z,\Gamma)$-affine.  Let $P_F^{T_0}\subset U$ be the minimal
canonically calibrated skeleton describing $T_0$ relative to $F$, as in
CLD Propositions~7.5.7--7.5.8.  In particular, $F$ is an immersion on each
cell of $P_F^{T_0}$; distinct analytic cells are not identified in this
description even if their images overlap in $F(U)$.  Choose a common
polyhedral subdivision $\mathscr C$ of $P_F^{T_0}$ on which all the
$\varphi_i\circ F$ are affine.  A top-dimensional cell $C$ carries its
canonical calibration
\[
                         \mu_C=m_C\mu_C^{\mathrm{prim}},
                         \qquad m_C\in\mathbf Z_{>0},
\]
where $\mu_C^{\mathrm{prim}}$ is the primitive lattice volume transported
through $F|_C$.  If $D$ is a codimension-one cell and $\nu_D$ is its
primitive lattice volume, the
formula in the proof of CLD Theorem~8.10.3 describes $T_i$ along $D$ by
\begin{equation}
\label{eq:p06r2-calibration}
                         \mu_{i,D}=a_{i,D}\nu_D,
                         \qquad a_{i,D}\in\mathbf Z.
\end{equation}
Indeed, the contribution from an adjacent $C$ is the canonical integer
weight $m_C$ times the derivative of the integral-affine function
$\varphi_i\circ F|_C$ in the primitive normal direction to $D$ in $C$, with the
orientation sign used in the boundary integral.  This derivative is an
integer; if $D$ is not primitive in the chosen ambient lattice, the usual
lattice index occurs, and it too is an integer.  Summing over the finitely
many adjacent cells gives $a_{i,D}\in\mathbf Z$.  This is exactly the
integral coefficient of the tropical Weil divisor of $\varphi_i\circ F$ on
the weighted skeleton.  A further common refinement subdivides these calibrated
cycles and preserves integrality.

By linearity and \eqref{eq:p06r2-independent}, the skeletal current
$\sum_i\lambda_iT_i=\mathrm d'\mathrm d''[h]$ is described on the same
subdivision along $D$ by
\begin{equation}
\label{eq:p06r2-sum-calibration}
                         \left(\sum_i\lambda_i a_{i,D}\right)\nu_D.
\end{equation}
The assumed current identity says that this skeletal current is zero.
Consequently the calibration in \eqref{eq:p06r2-sum-calibration} vanishes
on every codimension-one cell.  To make this last implication explicit,
if the calibration were nonzero on the relative interior of some $D$, one
could choose a constant Lagerberg $(n-1,n-1)$-form nonzero on $\nu_D$ and
a smooth cutoff supported in that relative interior.  The defining formula
of CLD Definition~7.5.1 would give a nonzero pairing.  Equivalently, this is
the uniqueness/minimality assertion of CLD Propositions~7.5.7--7.5.8.
Thus
\[
                         \sum_i\lambda_i a_{i,D}=0.
\]
Since $a_{i,D}\in\mathbf Z$ and the $\lambda_i$ are
$\mathbf Q$-linearly independent, $a_{i,D}=0$ for every $i$ and $D$.
Hence the calibrated skeletal description of each $T_i$ is zero in every
chart which tropicalizes the tuple $(h_1,\ldots,h_s)$.

We check that this proves the full current identity, not only a fixed-chart
pushforward.  Let $\alpha$ be any smooth compactly supported
$(n-1,n-1)$-test form in $W$.  By a smooth partition of unity and
Remark~7.2.5, it is a finite sum of terms $\rho\omega$, where $\rho$ is a
smooth compactly supported function and $\omega$ is tropicalized by a
moment map $H$ on an open set containing $\operatorname{Supp}(\rho)$.
Replace $F$ above by the common refinement $(F,H)$.  The same computation
on a common subdivision again has integral coefficients and shows that the
calibrated skeleton describing $T_i$ relative to $(F,H)$ is zero.  The
defining identity in CLD Definition~7.5.1 then gives
\[
                         \langle T_i,\rho\omega\rangle=0.
\]
Summing the localized terms proves
\eqref{eq:p06r2-integral-currents}.  Notice that no analytic current is
pulled back and no target representative is required to avoid the image of
the boundary of a compact analytic domain.  Compact support is carried by
$\rho$ on the open boundaryless space, exactly as in the definition of a
skeletal current.

For each $i$, apply CLD Theorem~18.8.4 with $d=1$ to the trivial line
bundle with the integral PL metric determined by $h_i$.  The separated and
boundaryless hypotheses are satisfied, and
\eqref{eq:p06r2-integral-currents} implies that $L_{h_i}(x)$ is numerically
trivial.  The residue construction is linear, so
\[
                         L_h(x)=\sum_i\lambda_iL_{h_i}(x)
\]
is numerically trivial.  Since $x$ was arbitrary, Gubler--Jell--Rabinoff,
Proposition~7.1 and Definition~7.2 imply that $h$ is harmonic.

Finally, harmonicity means that $h$ and $-h$ are semipositive.  By
Gubler--Rabinoff, Theorem~4.3, both functions are classically psh.  Hence
$h$ is classically pluriharmonic.
\end{proof}

\paragraph{Scope and rejection firewall.}
The proof applies only to $\mathbf R$-PL functions.  It does not pull back
arbitrary currents, does not use curve restriction, and does not assert the
boundary-separation statement rejected in production round~1.  It gives no
claim for arbitrary continuous functions.
